%! Radicals and Rational Exponents — Unit 1, Lesson 1.2 — Group Activity
\documentclass[10pt]{article}
\usepackage{atda-article}
\usepackage{atda-boxes}
\newcommand{\TallMath}[1]{$\displaystyle #1\rule[-1.4em]{0pt}{3.2em}$}
\setlength{\workrowsep}{6pt}

\begin{document}

\pageheader{Unit 1, Lesson 1.2}{Group Activity --- The Drop Tower}
% NAMESTRIP: no \namedateperiod here. The student writes their name once, on the
% cover this component is stapled behind. See references/conventions.md.

\vspace{0.15in}

\begin{scenariobox}[The Scenario]{royal}
\small
An amusement park is designing a free-fall drop tower. Ignoring air resistance, an object dropped
from rest falls $h=16t^2$ feet in $t$ seconds --- so the length of the free fall depends on the
height of the tower.

Work with your group. \textbf{Everyone starts at Tier R.} When your whole group can do Tier R
without help, move on.
\end{scenariobox}

\vspace{0.10in}

\boxguard[16]
\begin{tcolorbox}[colback=white, colframe=black!40, title=\textbf{Tier R --- Remediate},
    fonttitle=\bfseries, arc=2mm, left=3mm, right=3mm, top=2mm, bottom=2mm, breakable]
\small
Five fluency moves. Every one of them shows up again in Tier A.

\begin{enumerate}[leftmargin=1.6em, itemsep=4pt, topsep=4pt]
  \item Simplify: $\sqrt{75}$
\begin{work}
  \sqrt{75} &= \sqrt{25 \cdot 3} \\
            &= 5\sqrt{3}
\end{work}

  \item Simplify: $\sqrt[3]{16}$
\begin{work}
  \sqrt[3]{16} &= \sqrt[3]{8 \cdot 2} \\
               &= 2\sqrt[3]{2}
\end{work}

  \item Simplify: $\sqrt{45x^4}$
\begin{work}
  \sqrt{45x^4} &= \sqrt{9x^4 \cdot 5} \\
               &= 3x^2\sqrt{5}
\end{work}

  \item Combine: $\sqrt{20}+\sqrt{5}$
\begin{work}
  \sqrt{20}+\sqrt{5} &= 2\sqrt{5}+\sqrt{5} \\
                     &= 3\sqrt{5}
\end{work}

  \item Rewrite with a rational exponent: $\sqrt[3]{n^7}$
\begin{work}
  \sqrt[3]{n^7} &= n^{7/3}
\end{work}
\end{enumerate}
\end{tcolorbox}

\vspace{0.10in}

\boxguard[16]
\begin{tcolorbox}[colback=white, colframe=black!40,
    title=\textbf{Tier A --- Approaching Proficiency},
    fonttitle=\bfseries, arc=2mm, left=3mm, right=3mm, top=2mm, bottom=2mm, breakable]
\small
Now build the ride's model. A tower of height $h$ feet gives a free fall of $t$ seconds.

\begin{enumerate}[leftmargin=1.6em, itemsep=4pt, topsep=4pt]
  \item Start from $h=16t^2$ and solve for $t$. Show that the model can be written
        $t=\dfrac{\sqrt{h}}{4}$.
\begin{work}
  h &= 16t^2 \\
  \dfrac{h}{16} &= t^2 \\
  t &= \sqrt{\dfrac{h}{16}} = \dfrac{\sqrt{h}}{4}
\end{work}

  \item The park's first design is $128$ feet tall. Find the fall time exactly, then round to the
        nearest hundredth of a second.
\begin{work}
  t &= \dfrac{\sqrt{128}}{4} \\
    &= \dfrac{8\sqrt{2}}{4} \\
    &= 2\sqrt{2} \approx 2.83 \text{ seconds}
\end{work}

  \item Marketing wants a fall that lasts \emph{twice} as long. How tall must the tower be? Find
        the height, then write one sentence about what that costs the park.
\begin{work}
  4\sqrt{2} &= \dfrac{\sqrt{h}}{4} \\
  16\sqrt{2} &= \sqrt{h} \\
  h &= \left(16\sqrt{2}\right)^2 = 512 \text{ feet}
\end{work}
\writeline

  \item A second park writes its model as $t=\dfrac{12}{\sqrt{6}}$ seconds for its own tower.
        Rationalize the denominator and give the exact value.
\begin{work}
  \dfrac{12}{\sqrt{6}} &= \dfrac{12}{\sqrt{6}}\cdot\dfrac{\sqrt{6}}{\sqrt{6}} \\
                       &= \dfrac{12\sqrt{6}}{6} \\
                       &= 2\sqrt{6} \approx 4.90 \text{ seconds}
\end{work}

  \item The ride's safety fence is a rectangle measuring $\left(\sqrt{5}+2\right)$ meters by
        $\left(\sqrt{5}-4\right)$ meters on the plans. Multiply the two and simplify.
\begin{work}
  \left(\sqrt{5}+2\right)\left(\sqrt{5}-4\right) &= 5-4\sqrt{5}+2\sqrt{5}-8 \\
                                                 &= -3-2\sqrt{5}
\end{work}
\end{enumerate}
\end{tcolorbox}

\vspace{0.10in}

\boxguard[30]
\begin{tcolorbox}[colback=white, colframe=black!40, title=\textbf{Tier E --- Extension},
    fonttitle=\bfseries, arc=2mm, left=3mm, right=3mm, top=2mm, bottom=2mm, breakable]
\small

\begin{enumerate}[leftmargin=1.6em, itemsep=4pt, topsep=4pt]
  \item \textbf{Critique the work.} A student turned in these two answers. Both are wrong.
        Write the correct value for each, then name the mistake.
        \[ \sqrt{9+16} = 3+4 = 7 \qquad\text{and}\qquad 16^{1/2} = 8 \]
\begin{work}
  \sqrt{9+16} &= \sqrt{25} = 5 \\
  16^{1/2} &= \sqrt{16} = 4
\end{work}
\writeline

  \item \textbf{A tougher denominator.} Rationalize $\dfrac{6}{\sqrt{5}-2}$ by multiplying top and
        bottom by $\sqrt{5}+2$. Why does that particular factor clear the radical?
\begin{work}
  \dfrac{6}{\sqrt{5}-2} &= \dfrac{6\left(\sqrt{5}+2\right)}{\left(\sqrt{5}-2\right)\left(\sqrt{5}+2\right)} \\
                        &= \dfrac{6\sqrt{5}+12}{5-4} \\
                        &= 6\sqrt{5}+12 \approx 25.42
\end{work}
\writeline

  \item \textbf{Justify.} A cube-shaped water tank of volume $V$ cubic feet has edge length
        $s=V^{1/3}$. A second tank holds \emph{eight times} the water. Show that its edge is only
        twice as long, and explain why ``eight times the water'' does not mean ``eight times as
        tall.''
\begin{work}
  (8V)^{1/3} &= 8^{1/3}V^{1/3} \\
             &= 2V^{1/3} = 2s
\end{work}
\writelines{2}
\end{enumerate}
\end{tcolorbox}

\end{document}
